% ============================================================
%  paper_extended.tex
%  Accuracy Evaluation of MediaPipe Pose Estimation
%  Using a Unity Environment — Extended 8-page Review Paper
%
%  Base settings inherited from: docs/paper/original_paper.md
%  Target journal: IEEE Access / Sensors (MDPI)
%
%  Overleaf upload structure (relative to this file):
%    source/
%    ├── academic_paper/   ← this file lives here
%    ├── graphs/           ← heatmaps, correlation matrices
%    └── picture/          ← camera layout, concept diagrams, photos
%
%  Missing (from original IEEJ paper — supply separately):
%    source/picture/Unity実験環境.png
%    source/picture/frame_0018.jpg
% ============================================================

\documentclass[twocolumn]{article}

% ---- Package loading (original_paper.md baseline) ----------
\usepackage{graphicx}
\usepackage{geometry}
\usepackage{amsmath}
\usepackage{booktabs}
\usepackage{float}
\usepackage{url}
\usepackage{titlesec}

% ---- Extended packages (added for 8-page version) ----------
\usepackage{subcaption}   % subfigure layout (Fig 4, Fig 5, Fig 7/8 pairs)
\usepackage{multirow}     % merged rows in tables
\usepackage{array}        % column formatting
\usepackage{siunitx}      % angle unit \ang{}
\usepackage{hyperref}     % clickable references
\usepackage{cite}         % IEEE-style citation compression [1]–[4]

% ---- Page geometry (inherited from original_paper.md) ------
\geometry{top=24mm,bottom=24mm,left=19mm,right=19mm}
\setlength{\parskip}{0.15em}
\setlength{\floatsep}{6pt}
\setlength{\textfloatsep}{8pt}
\titlespacing*{\section}   {0pt}{2ex plus 0.5ex minus 0.2ex}{1.2ex plus 0.2ex}
\titlespacing*{\subsection}{0pt}{1.8ex plus 0.5ex minus 0.2ex}{0.9ex plus 0.2ex}
\titlespacing*{\subsubsection}{0pt}{1.5ex plus 0.3ex minus 0.2ex}{0.6ex plus 0.2ex}

\sloppy

% ---- Title block -------------------------------------------
\title{Accuracy Evaluation of MediaPipe Pose Estimation\\
       Using a Unity Simulation Environment}

\author{%
  Non-member\quad Fumima Taira$^{*a)}$ \qquad
  Member\quad    Tsukasa Kato$^{*b)}$ \qquad
  Non-member\quad Jin Afuso$^{*c)}$\footnotemark[1]
}
\date{}

% ============================================================
\begin{document}
\maketitle

\begingroup
\renewcommand{\thefootnote}{}
\footnotetext[1]{%
  $^{*a)}$ Okinawa Prefectural Kaiko High School,\quad
  $^{*b)}$ Graduate School of Education, University of the Ryukyus,\quad
  $^{*c)}$ lollol Inc.}
\endgroup

% ============================================================
% ABSTRACT  (target: ~0.5 page)
% ============================================================
\section*{Abstract}

%%TODO%% — Expand to 4–5 sentences. Include:
%  - Background: MediaPipe accuracy, viewpoint limitation of prior work
%  - Method: 505 cameras, 107 frames, 259,356 observations, 12 joints
%  - Two error types: Joint Angle MAE + Direction Angle Error (Δθ, Δψ)
%  - Key results (quantitative): elbow Δθ ~58°, hip Δψ r=−0.84, shoulder improvement after Y-flip
%  - Contributions as bullet list (3–4 items)

This study presents a large-scale quantitative evaluation of MediaPipe pose
estimation accuracy using Ground Truth collected in a Unity simulation
environment. Prior evaluations have largely been limited to frontal or
restricted camera configurations; evaluation across diverse viewpoints
essential in practical deployments has been insufficient.

We collected walking data from 505 multi-view cameras over 107 frames,
yielding 259,356 joint observations across 12 joints. Two error types were
evaluated: (1)~joint angle error (MAE) and (2)~direction angle error
($\Delta\theta$, $\Delta\psi$), together with inter-joint correlation
analysis. Our main findings are as follows.
\begin{itemize}\itemsep0pt
  \item Unifying the coordinate system (Y-axis orientation) reduces the mean
        elbow direction angle error from $\approx121°$ to $\approx58°$, and
        shoulder error from $\approx120°$ to $\approx10°$.
  \item Left and right hip depth errors exhibit a strong negative correlation
        ($r = -0.84$), revealing a systematic pelvic rotation artefact.
  \item Upper-limb errors are strongly coupled ($r = 0.72$–$0.77$),
        suggesting hierarchical error propagation.
  \item Error patterns are viewpoint-dependent and largely predictable,
        offering scope for post-processing correction.
\end{itemize}

\vspace{-0.8em}

% ============================================================
% 1. INTRODUCTION  (target: ~1.5 pages)
% ============================================================
\section{Introduction}

Pose estimation has become a core enabling technology across sports
analysis, clinical rehabilitation, ergonomics, and human–computer
interaction~\cite{mediapipe,blazepose}. MediaPipe, developed by Google, is
particularly attractive for practitioners because it enables real-time
processing on commodity hardware without GPU acceleration~\cite{mediapipe}.

Despite its widespread deployment, a systematic and multi-viewpoint
quantitative evaluation of MediaPipe's 3D accuracy remains lacking.
Bazarevsky~et~al.~\cite{blazepose} report BlazePose benchmark performance,
but their evaluation is restricted to a curated benchmark with limited
viewpoint diversity. In real-world deployments, camera placement is
constrained by the environment and cannot be controlled; understanding how
accuracy varies with viewpoint is therefore essential.

%%TODO%% — Expand this section with:
%  (a) Prior work on pose estimation accuracy evaluation (see Related Work
%      section — cross-reference here briefly)
%  (b) The monocular depth ambiguity problem
%  (c) Role of simulation/Unity-based GT collection
%  (d) Gap this paper fills

The contributions of this paper are as follows:
\begin{enumerate}\itemsep0pt
  \item \textbf{Large-scale multi-viewpoint evaluation}: 505 camera
        positions spanning four height layers and a wide horizontal range,
        yielding 259,356 observations — the largest multi-viewpoint
        MediaPipe evaluation to date.
  \item \textbf{Coordinate system unification}: We identify and correct a
        systematic Y-axis orientation mismatch between Unity and MediaPipe,
        and demonstrate its substantial effect on elbow direction accuracy.
  \item \textbf{Direction angle error analysis}: We introduce and quantify
        $\Delta\theta$ and $\Delta\psi$ as complementary metrics to joint
        angle MAE, revealing bias patterns invisible to MAE alone.
  \item \textbf{Inter-joint error correlation}: Pearson correlation analysis
        across 12 joints exposes systematic model-level constraints including
        pelvic anti-correlation and upper-limb coupling.
\end{enumerate}

% ============================================================
% 2. RELATED WORK  (target: ~1 page) — NEW SECTION
% ============================================================
\section{Related Work}

%%TODO%% — Write approximately 1 full page covering:

\subsection{MediaPipe and BlazePose Accuracy}
%%TODO%% — Summarise Bazarevsky et al. 2020 \cite{blazepose}.
%  Report their benchmark numbers (PCK on COCO/LSP).
%  Note the limitation: frontal-biased viewpoints.
%  Cite 1–2 independent evaluations of MediaPipe on real humans.

\subsection{Simulation-Based Ground Truth Collection}
%%TODO%% — Describe precedents of using game engines / Unity to obtain GT:
%  - Synthetic data (e.g., SURREAL, AGORA) for pose estimation training
%  - Unity-based capture systems for biomechanics evaluation
%  - Advantages: perfect GT, controllable conditions, viewpoint sampling

\subsection{Monocular 3D Pose Estimation and Depth Ambiguity}
%%TODO%% — Discuss:
%  - Inherent depth ambiguity in monocular lifting (Mehta et al.)
%  - How training dataset viewpoint bias (Human3.6M \cite{h36m}, COCO \cite{coco})
%    affects generalisation to non-frontal viewpoints
%  - Limitations of skeleton models without rigidity constraints (pelvis as line)

\subsection{Skeletal Constraint Models}
%%TODO%% — GNN-based skeleton models, bone-length constraints, pelvis rigidity.
%  Reference 1–2 recent papers (e.g., PoseFormer, MotionBERT).

% ============================================================
% 3. EXPERIMENTAL METHOD  (target: ~1.5 pages)
% ============================================================
\section{Experimental Method}

\subsection{Experimental Environment}

Using Unity (version 6000.0.60f1), we constructed a walking simulation
environment with a Y-Bot humanoid avatar. The avatar performed a straight
walking motion from coordinates $(0,\,0,\,-3)$ to $(0,\,0,\,3)$ along the
positive Z-axis.

\begin{figure}[!t]
\centering
\includegraphics[width=0.30\textwidth]{figs/fig01_camera_layout.png}
\caption{Camera layout (505 positions across four height layers Y = 0.5,
  1.0, 1.5, 2.0). The centre arrow indicates the avatar's walking direction
  (+Z axis).}
\label{fig:camera_layout}
\end{figure}

Cameras were placed on a 3-D grid covering four height layers
(Y = 0.5, 1.0, 1.5, 2.0\,m) and horizontal positions
(X, Z $\in [-6, 6]$\,m). Within each layer a $13\times13$ grid was used
with the central $5\times5$ block removed, giving a theoretical maximum of
$4(13^2-5^2)=576$ positions. Data were collected at 505 of these positions.
Each position yielded 107 captured frames at a resolution of
$1280\times720$\,pixels, for a total of 259,356 joint observations
(107 frames $\times$ 505 cameras $\times$ 12 joints, after visibility
filtering). The same GT and image data are reusable for evaluating
alternative estimators such as YOLO-Pose.

%%TODO%% — Add 1–2 sentences on avatar parameters (Y-Bot, animation clip name,
%  walking speed/cycle if known).

% NOTE (NEW_PAPER_MANAGEMENT):
%   Images/Unity実験環境.png および Images/frame_0018.jpg は
%   元の投稿論文（docs/paper/Zeval_IEEJ_submit.pdf）に掲載済みの図。
%   Zeval_DataSet_organized/07_paper/Images/ への配置が未完了。
%   元論文から抽出するか、Unity/1_Output_Photos から取得すること。
\begin{figure}[!t]
  \centering
  \includegraphics[width=0.44\textwidth]{figs/fig02a_unity_env.png}
  \vspace{0.3em}
  \includegraphics[width=0.44\textwidth]{figs/fig02b_frame_sample.jpg}
  \caption{Unity experiment environment (top) and a representative captured
    frame (bottom).}
  \label{fig:unity_env}
\end{figure}

\subsection{Capture System}

A Trigger Zone Capture System was implemented in Unity.
\texttt{CaptureSystemManager} handled character generation, animation
playback, per-frame image capture, and joint coordinate recording.
\texttt{AutoCaptureManager} read the 505 camera positions from a CSV file
and triggered automatic capture at each position. This ensured consistent,
reproducible multi-viewpoint data collection with a single character motion
sequence.

\subsection{Data Collection}

Joint coordinates (GT) were exported per frame in CSV format using Unity's
\texttt{HumanBodyBones} API~\cite{unity}. MediaPipe's Pose Landmarker
(\texttt{model\_complexity=1}, \texttt{static\_image\_mode=True},
\texttt{min\_detection\_confidence=0.5}) was applied to each saved image to
extract 3-D coordinates for 33 landmarks. Landmarks with a visibility score
below 0.5 were excluded from analysis to limit the effect of heavy
occlusion. The 12 joints retained for analysis are: left/right shoulder,
elbow, wrist, hip, knee, and ankle.

\subsection{Angle Computation and Error Evaluation}

We report \textbf{two complementary error metrics}.

\subsubsection{Joint Angle Error (Joint Angle MAE)}

The flexion angle at each joint was computed from three anatomical points:
upper arm–elbow–forearm for the elbow; thigh–knee–lower leg for the knee;
and equivalent triplets for shoulder and hip. The angle $\alpha$ at the
vertex between the two vectors was computed as
\[
  \alpha = \arccos\!\left(\frac{\mathbf{u}\cdot\mathbf{v}}
                               {\|\mathbf{u}\|\,\|\mathbf{v}\|}\right),
  \quad \alpha \in [0°,\,180°].
\]
The mean absolute difference between GT and MediaPipe angles over all valid
frames and cameras gives the MAE for each joint.

\subsubsection{Direction Angle Error ($\Delta\theta$, $\Delta\psi$)}

Following the BlazePose output specification~\cite{blazepose}, relative
joint positions are expressed with the hip midpoint as origin. In this
hip-centred coordinate system we define
\[
  \theta = \mathrm{arctan2}(y,\,x), \quad
  \psi   = \mathrm{arctan2}(z,\,x),
\]
where $(x,y,z)$ is the joint position relative to the hip centre.
The angular errors $\Delta\theta$ and $\Delta\psi$ are the signed
differences between GT and MediaPipe values, normalised to
$(-180°, +180°]$.

\subsubsection{Coordinate System Unification}

Unity uses a left-handed coordinate system with Y-axis pointing
upward, whereas MediaPipe's \texttt{pose\_landmarks} uses an image
coordinate system with Y-axis pointing downward. This discrepancy
introduces a systematic sign flip in $\theta = \mathrm{arctan2}(y,x)$.
Before computing direction angle errors we applied the transformation
\[
  y_{\text{MediaPipe}}^{\text{unified}} = -y_{\text{MediaPipe}},
\]
implemented in
\texttt{06\_direction\_detection/scripts/coordinate\_transform.py}.
All direction angle results and inter-joint correlation analyses reported
below use coordinates after this unification.

\begin{figure}[!t]
\centering
\includegraphics[width=0.48\linewidth]{figs/fig03_concept_metrics.png}
\caption{Concept diagram of the two evaluation metrics. Left: joint angle
  (MAE) — flexion angle of three points ($0$–$180°$). Right: direction
  angle ($\Delta\theta$, $\Delta\psi$) — angular offset relative to the
  hip centre ($-180°$ to $+180°$). The figure is a simplified schematic;
  the actual analysis uses all 33 MediaPipe landmarks~\cite{blazepose}.}
\label{fig:concept_metrics}
\end{figure}

% ============================================================
% 4. RESULTS  (target: ~2.5 pages)
% ============================================================
\section{Results}

\subsection{Joint Angle MAE}

Table~\ref{tab:joint_angle_error} summarises MAE, median, and maximum
error for eight major joints. Elbow and knee joints exhibit the lowest
errors (~18° and ~17--19°, respectively), while the shoulder joints show
the highest MAE (~39--40°), likely reflecting the increased depth ambiguity
of proximal joints relative to the camera.

\begin{table}[!t]
\centering
\caption{Statistics of joint angle error (Joint Angle MAE, degrees).
  Source: \texttt{04\_mae\_heatmap/coordinate\_angle\_mae.csv}.}
\label{tab:joint_angle_error}
\begin{tabular}{lccc}
\toprule
Joint          & MAE    & Median & Max    \\
               & (deg)  & (deg)  & (deg)  \\
\midrule
Left shoulder  & 39.8   & 39.2   & 70.1   \\
Right shoulder & 39.0   & 40.0   & 58.2   \\
Left elbow     & 18.6   & 15.9   & 56.9   \\
Right elbow    & 18.2   & 13.6   & 46.2   \\
Left hip       & 30.2   & 29.4   & 66.0   \\
Right hip      & 33.3   & 32.2   & 81.7   \\
Left knee      & 17.0   & 16.6   & 43.4   \\
Right knee     & 19.5   & 19.4   & 62.3   \\
\bottomrule
\end{tabular}
\end{table}

%%TODO%% — Add Y-layer comparison sub-table or figure:
%  Y=0.5,1.5 vs Y=1.0,2.0 split of MAE per joint.
%  Source: 04_mae_heatmap/Y=0.5,1.5/ and Y=1.0.2.0/ outputs.
%  Data to be filled in by data-collection agent — see DATA_COLLECTION_AGENT_INSTRUCTIONS.md §2.

\subsection{Direction Angle Error ($\Delta\theta$, $\Delta\psi$)}

Table~\ref{tab:direction_angle_error} presents the direction angle error
statistics for all 12 joints after coordinate system unification.

\begin{table}[!t]
\centering
\caption{Direction angle error statistics after coordinate unification
  (degrees; absolute-value means). Source:
  \texttt{11\_direction\_ditection/output/processed\_data/joint\_summary.csv}.
  $^\dagger$Hip $|\Delta\psi|$ from separate analysis (see text).}
\label{tab:direction_angle_error}
\small
\begin{tabular}{lcccc}
\toprule
Joint & \multicolumn{2}{c}{$|\Delta\theta|$} & \multicolumn{2}{c}{$|\Delta\psi|$} \\
\cmidrule(lr){2-3}\cmidrule(lr){4-5}
      & Mean & SD & Mean & SD \\
\midrule
Left shoulder  & 10.4 & 11.7 & 94.1 & 86.4 \\
Right shoulder &  9.3 & 10.1 & 92.7 & 92.0 \\
Left elbow     & 58.3 & 24.1 & 87.3 & 90.6 \\
Right elbow    & 59.3 & 22.7 & 88.6 & 94.5 \\
Left wrist     & 84.3 & 103.7 & 88.6 & 91.9 \\
Right wrist    & 91.5 & 111.9 & 91.6 & 99.2 \\
Left hip$^\dagger$ & — & — & 88.9 & — \\
Right hip$^\dagger$& — & — & 91.1 & — \\
Left knee      & 17.4 & 18.8 & 84.7 & 97.9 \\
Right knee     & 18.5 & 19.1 & 87.2 & 97.5 \\
Left ankle     & 11.3 & 14.0 & 90.6 & 105.8 \\
Right ankle    & 11.9 & 14.5 & 100.4 & 103.1 \\
\bottomrule
\end{tabular}
\end{table}

Figure~\ref{fig:elbow_heatmaps} shows the camera-position-wise distribution
of elbow $\Delta\theta$.

\begin{figure}[!t]
  \centering
  \begin{subfigure}[b]{0.48\linewidth}
    \includegraphics[width=\linewidth]%
      {figs/fig04a_elbow_L_theta_y05.jpg}
    \caption{Left elbow $\Delta\theta$ (Y=0.5)}
  \end{subfigure}\hfill
  \begin{subfigure}[b]{0.48\linewidth}
    \includegraphics[width=\linewidth]%
      {figs/fig04b_elbow_R_theta_y05.jpg}
    \caption{Right elbow $\Delta\theta$ (Y=0.5)}
  \end{subfigure}
  \caption{Heatmaps of elbow direction angle error $\Delta\theta$ (Y=0.5).
    Left and right elbows show opposite-sign patterns, consistent with the
    strong negative inter-joint correlation ($r=-0.862$).}
  \label{fig:elbow_heatmaps}
\end{figure}

Figure~\ref{fig:hip_psi_heatmap} shows the hip $\Delta\psi$ heatmap.
Left and right hip exhibit near-identical absolute-value patterns due to
their strong negative correlation ($r=-0.84$).

\begin{figure}[!t]
  \centering
  \includegraphics[width=0.75\linewidth]%
    {figs/fig05a_hip_L_psi_y05.jpg}
  \caption{Heatmap of hip joint direction angle error $\Delta\psi$
    (Y=0.5, left hip). The right hip shows an almost identical
    absolute-value pattern owing to the anti-correlation
    $\Delta\psi_{\text{L}} \approx -\Delta\psi_{\text{R}}$.}
  \label{fig:hip_psi_heatmap}
\end{figure}

Table~\ref{tab:mae_by_layer} breaks down joint angle MAE by camera height
layer. Errors are generally lowest at Y=1.0 and increase at Y=2.0
(overhead view), suggesting that top-down perspectives reduce estimation
accuracy, consistent with underrepresentation of such viewpoints in training
data.

\begin{table}[!t]
\centering
\caption{Joint angle MAE (deg) by camera height layer.
  Source: \texttt{4\_MAE\_HEATMAP/Y=0.5,1.5/} and
  \texttt{4\_MAE\_HEATMAP/Y=1.0.2.0/coordinate\_angle\_mae.csv}.}
\label{tab:mae_by_layer}
\small
\begin{tabular}{lcccc}
\toprule
Joint & Y=0.5 & Y=1.0 & Y=1.5 & Y=2.0 \\
\midrule
Left shoulder  & 39.8 & 38.9 & 39.6 & 41.0 \\
Right shoulder & 39.3 & 38.1 & 38.6 & 40.2 \\
Left elbow     & 18.1 & 17.7 & 18.3 & 20.2 \\
Right elbow    & 17.5 & 16.7 & 18.2 & 20.2 \\
Left hip       & 29.6 & 29.0 & 29.9 & 32.5 \\
Right hip      & 32.0 & 31.7 & 32.8 & 36.7 \\
Left knee      & 16.1 & 16.6 & 17.7 & 17.8 \\
Right knee     & 18.9 & 18.7 & 19.6 & 20.6 \\
\bottomrule
\end{tabular}
\end{table}

\subsection{Inter-joint Error Correlation}

Tables~\ref{tab:correlation_theta} and \ref{tab:correlation_psi} list
high-correlation pairs ($|r|>0.7$) for $\Delta\theta$ and $\Delta\psi$,
respectively.

\begin{table}[!t]
\centering
\caption{High-correlation joint pairs in XY plane ($\Delta\theta$).}
\label{tab:correlation_theta}
\small
\begin{tabular}{llr}
\toprule
Joint 1        & Joint 2        & Correlation \\
\midrule
LEFT\_ELBOW    & RIGHT\_ELBOW   & $-0.862$ \\
LEFT\_ELBOW    & LEFT\_SHOULDER &  $0.788$ \\
RIGHT\_ANKLE   & RIGHT\_KNEE    &  $0.767$ \\
LEFT\_ANKLE    & LEFT\_KNEE     &  $0.766$ \\
RIGHT\_ELBOW   & RIGHT\_SHOULDER&  $0.710$ \\
\bottomrule
\end{tabular}
\end{table}

\begin{table}[!t]
\centering
\caption{High-correlation joint pairs in XZ plane ($\Delta\psi$).}
\label{tab:correlation_psi}
\small
\begin{tabular}{llr}
\toprule
Joint 1        & Joint 2         & Correlation \\
\midrule
LEFT\_HIP      & RIGHT\_HIP      & $-0.840$ \\
RIGHT\_ELBOW   & RIGHT\_SHOULDER &  $0.770$ \\
RIGHT\_ELBOW   & RIGHT\_WRIST    &  $0.769$ \\
LEFT\_ELBOW    & LEFT\_SHOULDER  &  $0.768$ \\
RIGHT\_SHOULDER& RIGHT\_WRIST    &  $0.726$ \\
LEFT\_ELBOW    & LEFT\_WRIST     &  $0.722$ \\
RIGHT\_ELBOW   & RIGHT\_HIP      &  $0.721$ \\
LEFT\_HIP      & RIGHT\_ELBOW    & $-0.705$ \\
\bottomrule
\end{tabular}
\end{table}

In the XY plane, the dominant feature is a strong negative correlation
between left and right elbow ($r=-0.862$), alongside positive same-side
elbow–shoulder and ankle–knee coupling. In the XZ plane, the hip
anti-correlation ($r=-0.84$) and upper-limb coupling ($r=0.72$–$0.77$) are
notable. Figure~\ref{fig:correlation_heatmaps} visualises the full
correlation matrices.

\begin{figure*}[t]
  \centering
  \begin{subfigure}[b]{0.49\textwidth}
    \centering
    \includegraphics[width=\linewidth]{figs/fig06a_corr_theta.png}
    \caption{$\Delta\theta$ (XY plane)}
  \end{subfigure}\hfill
  \begin{subfigure}[b]{0.49\textwidth}
    \centering
    \includegraphics[width=\linewidth]{figs/fig06b_corr_psi.png}
    \caption{$\Delta\psi$ (XZ plane)}
  \end{subfigure}
  \caption{Correlation matrices of inter-joint direction angle errors.
    (a)~$\Delta\theta$ corresponds to Table~\ref{tab:correlation_theta};
    (b)~$\Delta\psi$ corresponds to Table~\ref{tab:correlation_psi}.
    Source: \texttt{06\_direction\_detection/output/correlation\_analysis/}.}
  \label{fig:correlation_heatmaps}
\end{figure*}

\subsection{Camera Viewpoint Dependency}

To characterise viewpoint dependence, we identified for each of the
107 frames the camera with the lowest and highest mean $|\Delta\theta|$
across joints (computed from
\texttt{11\_direction\_ditection/output/processed\_data/frame\_camera\_summary.csv},
107 frames $\times$ 201 cameras per frame).

The \emph{best} camera per frame achieves a mean $|\Delta\theta|$ of
$8.8°$ averaged over all frames (min $6.1°$, max $15.8°$), demonstrating
that near-optimal viewpoints exist. In contrast, the \emph{worst} camera
yields a mean of $80.1°$ (min $55.6°$, max $148.1°$), indicating that
highly unfavourable viewpoints produce errors comparable to random
direction estimates. This 9-fold difference between best and worst cameras
underscores the critical importance of camera placement in practical
deployments.

% ============================================================
% 5. DISCUSSION  (target: ~2 pages)
% ============================================================
\section{Discussion}

\subsection{Coordinate System Unification Effect}

The Y-axis orientation mismatch between Unity and MediaPipe causes a
systematic sign error in $\theta$ for every joint. Because
$\theta_{\text{MP}} \approx -\theta_{\text{GT}}$ when the Y-axis is
inverted, the uncorrected direction angle error satisfies
$\Delta\theta \approx -2\theta_{\text{GT}}$; for a joint at
$\theta_{\text{GT}} \approx 60°$ this yields $|\Delta\theta| \approx 120°$.
After applying $y \leftarrow -y$ to MediaPipe coordinates, the mean
$|\Delta\theta|$ for \emph{shoulder} joints dropped from $\approx120°$ to
$\approx10°$, and for \emph{elbow} joints from $\approx121°$ to
$\approx58°$~\cite{mediapipe_coord_issue}.

% NOTE: The "before Y-flip" values (~121° for elbow) are documented in
% 10_theta_verification/README.md and MEDIAPIPE_COORDINATE_SYSTEM.md.
% Precise per-joint before/after table requires re-running
% 10_theta_verification/test_04_coordinate_system.py without the Y-flip.

The residual $\approx58°$ elbow error after unification reflects genuine
estimation bias (training-data viewpoint imbalance, arm-at-side prior),
not coordinate system mismatch.

\subsection{Systematic Bias in Depth Estimation}

MediaPipe's BlazePose estimates each landmark via heatmaps and regression,
outputting 3-D world coordinates relative to the hip
centre~\cite{blazepose}. The architecture does not explicitly incorporate
anatomical constraints such as bone-length consistency or pelvic
rigidity. The mean $|\Delta\theta|$ for elbow joints after coordinate
unification remains $\approx58°$, indicating a persistent directional bias.

Many pose estimation models, including MediaPipe, are trained on public
datasets such as Human3.6M~\cite{h36m} and COCO~\cite{coco} in which
everyday, near-frontal, standing/walking poses dominate. The arm-at-side
prior embedded in these datasets likely biases the elbow direction estimate
toward a downward-hanging posture. The opposite-sign pattern in left vs.\
right elbow errors ($r=-0.862$) strongly suggests left–right flip data
augmentation during training; the standard deviation of elbow $\Delta\theta$
($\approx23$–$25°$) is small enough that the bias is relatively
predictable and correctable by post-processing.

\subsection{Left-Right Symmetric Hip Depth Error}

For hip joints, the mean $|\Delta\psi|\approx90°$ and the strong negative
correlation ($r=-0.8402$) indicate that the depth-direction errors are
anti-correlated: when one hip is estimated as forward, the other is
estimated as backward, creating a phantom pelvic-rotation artefact. The
identical absolute-value heatmap patterns for left and right hip
(Figure~\ref{fig:hip_psi_heatmap}) are consistent with
$\Delta\psi_{\text{L}} \approx -\Delta\psi_{\text{R}}$ at each observation.

In MediaPipe's skeleton model the pelvis is represented as the segment
connecting landmark~23 (left hip) and landmark~24 (right hip). This linear
representation without a multi-segment rigid structure (e.g., sacrum and
ilium) imposes no rigidity constraint on the relative depth of the two
hips during 3-D lifting. Consequently, even a small horizontal pixel
difference in the 2-D image can be interpreted as placing one hip closer
and the other farther from the camera, producing the symmetric error
pattern. Model improvements that impose pelvis rigidity as a hard or soft
constraint are needed.

\subsection{Upper-Limb Coupled Error}

The high positive correlations among shoulder, elbow, and wrist in
$\Delta\psi$ ($r=0.72$–$0.77$) suggest that errors propagate
hierarchically along the kinematic chain, or that the upper limb is
implicitly treated as a rigid unit during 3-D lifting. Incorporating
anatomical bone-length constraints via graph neural networks or explicit
kinematic models could reduce this propagation.

\subsection{Camera Viewpoint Dependency}

Selecting the best camera per frame reduces $|\Delta\theta|$ to 6.1°–15.8°
(mean 8.8°), while the worst camera yields 55.6°–148.1° (mean 80.1°).
This 9-fold range between best and worst viewpoints for the same body pose
confirms that estimation quality is heavily viewpoint-dependent.
The training data for BlazePose (primarily Human3.6M~\cite{h36m} and
COCO~\cite{coco}) is dominated by frontal and near-frontal viewpoints;
cameras at extreme angles (lateral, top-down) are underrepresented, leading
to reduced generalisation. This is consistent with the elevated MAE at
Y=2.0 (overhead cameras) observed in Table~\ref{tab:mae_by_layer}.

\subsection{Practical Implications}

MediaPipe provides useful 2-D joint localisation but its 3-D depth
estimation is fundamentally limited by monocular ambiguity.
For applications where relative joint angles in the image plane suffice
(e.g., stroke-count analysis in swimming, gross-motor rehabilitation
screening) MediaPipe is adequate. For applications requiring 3-D body
orientation (e.g., gait analysis, ergonomic risk assessment of torso
rotation) the systematic biases identified here—particularly in hip
$\Delta\psi$ and elbow $\Delta\theta$—must be compensated.

\subsection{Limitations}

This study has the following limitations. (1)~The avatar is a standard
Y-Bot mesh; skin deformation and clothing effects present in real subjects
are absent. (2)~Only a single walking motion was evaluated; other motion
types may produce different error patterns. (3)~The simulated lighting is
uniform; real-world illumination variation was not assessed. (4)~Only
MediaPipe was evaluated; comparison with other estimators (e.g., YOLO-Pose,
OpenPose) is left for future work.

% ============================================================
% 6. CONCLUSION  (target: ~0.5 page)
% ============================================================
\section{Conclusion}

We presented a large-scale multi-viewpoint evaluation of MediaPipe pose
estimation using a Unity simulation environment. The key findings are:
\begin{itemize}\itemsep0pt
  \item Joint angle MAE ranges from $\approx17°$ (knee) to $\approx40°$
        (shoulder); elbow and knee errors are the lowest.
        MAE increases at Y=2.0 (overhead) cameras, consistent with
        training-data viewpoint bias.
  \item Coordinate system unification (Y-axis flip) is essential:
        it reduces shoulder $|\Delta\theta|$ from $\approx120°$ to
        $\approx10°$ and elbow $|\Delta\theta|$ from $\approx121°$ to
        $\approx58°$.
  \item Elbow direction error ($|\Delta\theta|\approx58°$) reflects
        training-data bias and is systematic (SD $\approx23$–$25°$),
        offering scope for post-processing correction.
  \item Hip depth error ($|\Delta\psi|\approx90°$, $r=-0.84$) reflects
        an absence of pelvic rigidity constraints in the skeleton model.
  \item Upper-limb errors are strongly coupled ($r=0.72$–$0.77$),
        suggesting hierarchical error propagation.
\end{itemize}

Future directions include: (1)~incorporating pelvic rigidity constraints
into pose estimation models; (2)~fine-tuning on multi-viewpoint data
collected from diverse angles; (3)~extending the evaluation framework to
other estimators and motion types; and (4)~developing application-specific
correction filters for elbow and hip errors.

% ============================================================
\section*{Acknowledgments}
% ============================================================

This work was supported by the Japan Science and Technology Agency (JST)
under the Next-Generation Human Resource Development Program (Global Science
Campus Initiative).

% ============================================================
% REFERENCES
% ============================================================
\begin{thebibliography}{99}

% --- Original 3 references (kept from paper_en.md) ----------
\bibitem{mediapipe}
Google LLC, ``MediaPipe Pose Landmarker,''
\url{https://ai.google.dev/edge/mediapipe/solutions/vision/pose_landmarker},
2024.

\bibitem{mediapipe_coord_issue}
Google, ``Orientation of pose world landmarks,''
GitHub Issue \#3370, MediaPipe repository, 2022.
\url{https://github.com/google/mediapipe/issues/3370}

\bibitem{blazepose}
V.~Bazarevsky, I.~Grishchenko, K.~Raveendran, T.~Zhu, F.~Zhang,
and M.~Grundmann,
``BlazePose: On-device real-time body pose tracking,''
\textit{arXiv preprint arXiv:2006.10204}, 2020.

\bibitem{unity}
Unity Technologies, ``HumanBodyBones,''
\textit{Unity Scripting Reference}, 2024,
\url{https://docs.unity3d.com/ScriptReference/HumanBodyBones.html}.

% --- New references to be added (%%TODO%% — verify BibTeX) --
\bibitem{h36m}
%%TODO%% C.~Ionescu et al.,
``Human3.6M: Large Scale Datasets and Predictive Methods for 3D Human
Sensing in Natural Environments,''
\textit{IEEE Trans. Pattern Anal. Mach. Intell.}, vol.~36, no.~7,
pp.~1325--1339, 2014.

\bibitem{coco}
%%TODO%% T.-Y.~Lin et al.,
``Microsoft COCO: Common Objects in Context,''
in \textit{Proc. ECCV}, 2014, pp.~740--755.

\bibitem{openpose}
%%TODO%% Z.~Cao, G.~Hidalgo, T.~Simon, S.-E.~Wei, and Y.~Sheikh,
``OpenPose: Realtime Multi-Person 2D Pose Estimation Using Part Affinity
Fields,''
\textit{IEEE Trans. Pattern Anal. Mach. Intell.}, vol.~43, no.~1,
pp.~172--186, 2021.

% %%TODO%% — Add 4–6 more references from Related Work section:
%   - 1–2 independent MediaPipe accuracy papers
%   - 1   Unity/synthetic GT dataset paper (e.g., SURREAL or AGORA)
%   - 1   monocular depth lifting paper (e.g., Mehta et al. 2017)
%   - 1   skeletal constraint / GNN paper (e.g., PoseFormer or MotionBERT)

\end{thebibliography}

\end{document}

